\section{A symmetry obstruction for a larger inducing field}
\label{sec:inducing-symmetry}

\begin{theorem}[D-BI-15: normalizer and unique-system obstruction]
\label{D-BI-15}
Let $C\to\mathbf P^1$ be a connected finite $G$-cover in a
Hurwitz family, and let $\rho$ be an absolutely irreducible
representation with a $K$-model and exact character field $K$.
Suppose $\rho=\operatorname{Ind}_H^G\tau$, put
$L=\mathbf Q(\chi_\tau)$, and assume $L\supsetneq K$.
Suppose the actual intermediate base $X=C/H$ is rational and
the faithful cover $C/\ker\tau\to X$ has $R\ge3$ genuine
branch points.  If there exist
$\gamma\in\operatorname{Gal}(\overline{\mathbf Q}/K)$ and
$g\in N_G(H)$ such that
\[
                 \tau^\gamma\not\simeq\tau,
                 \qquad \tau^\gamma\simeq\tau^g,
                                                               \tag{BI.51}
\]
then the period rank of its multiplicity variation is at most
\[
                         \lfloor R/2\rfloor-1.             \tag{BI.52}
\]
In particular a twelve-branch intermediate cover of this kind
cannot supply period rank nine.

Condition (BI.51) follows in either of the following two cases:
$H$ is normal in $G$; or the given system of imprimitivity of
$\rho$ is the unique system with its block dimension.  In the
second assertion uniqueness among primitive systems also suffices,
if the inducing representation is primitive.  No general
uniqueness theorem for imprimitivity systems is assumed.
\end{theorem}
\begin{proof}
Write $N=\ker\tau$.  Galois conjugation preserves this kernel.
Indeed a finite-order matrix has character equal to its dimension
exactly when it is the identity: its eigenvalues have modulus one,
and equality in the triangle inequality forces them all to equal
one.  The equality $\chi_\tau(h)=\dim\tau$ is preserved and
reflected by $\gamma$.  Thus (BI.51) implies $gNg^{-1}=N$.
Moreover $g\notin H$, since inner conjugation leaves the
isomorphism class of $\tau$ unchanged.

The element $g$ consequently induces a nontrivial automorphism
$\varphi$ of $X=C/H$, and lifts to an automorphism of $C/N$.
The action of $N_G(H)/H$ on $C/H$ is faithful, by the
Galois correspondence for the connected $G$-cover.  Hence
$\varphi$ has finite order $q\ge2$ and preserves the reduced
branch divisor of $C/N\to X$.  This statement concerns exactly
the genuine branch points, since an isomorphism of covers
preserves ramification.

Here is the dimension bound without any assumption about a
generic automorphism group.  Conjugate a finite-order automorphism
of $\mathbf P^1$ to $z\mapsto\zeta_qz$.  An invariant reduced
divisor of degree $R$ consists of $f\in\{0,1,2\}$ of the fixed
points $0,\infty$ and $s=(R-f)/q$ nonfixed orbits.  Since
$R\ge3$, one has $s\ge1$ and $q\le R$.  The $s$ orbit
parameters lie in $\mathbf C^*$; quotienting by the one-dimensional
centralizer of the automorphism leaves dimension $s-1$.
Choices of labels and of the fixed points are finite.  There are
only finitely many possible $q$ and $f$, and therefore the locus
of such branch configurations in $M_{0,R}$ has dimension at most
\[
             s-1\le\lfloor R/2\rfloor-1.
\]
For fixed finite monodromy, the intermediate Hurwitz map to its
marked-base moduli is finite over the smooth locus.  Frobenius
reciprocity identifies the original multiplicity variation with
the one on this intermediate cover.  Its period rank is thus
bounded by the dimension just computed.  This proves (BI.52).

It remains to justify the two sufficient conditions.  Choose
$\gamma$ fixing $K$ but nontrivial on $L$, so that
$\tau^\gamma\not\simeq\tau$.  If $H\triangleleft G$, then
\[
       \operatorname{Res}_H^G\operatorname{Ind}_H^G\tau
                      =\bigoplus_{G/H}\tau^g.
\]
Since $\rho^\gamma\simeq\rho$, the constituent $\tau^\gamma$
must be one of these conjugates.  Normality puts the corresponding
$g$ in $N_G(H)$, proving (BI.51).

For the uniqueness assertion, let $W$ be the chosen inducing
block, whose stabilizer is $H$.  Frobenius reciprocity gives
\[
 \dim\operatorname{Hom}_H(\tau,\operatorname{Res}_H^G\rho)
 =\dim\operatorname{Hom}_G(\operatorname{Ind}_H^G\tau,\rho)=1.
\]
Thus $W$ is the image of the central $\tau$-projector in
$\overline{\mathbf Q}[H]$ acting on the $K$-model of $\rho$;
in particular the subspace and its translates are defined over
$\overline{\mathbf Q}$.  Applying $\gamma$ to this system gives
another $G$-stable system of the same block dimension, because
all matrices of $\rho$ are fixed by $\gamma$.  If the original
inducing representation is primitive, so is its Galois conjugate:
induction from a proper subgroup commutes with Galois conjugation,
and the existence of an imprimitivity system is equivalent to
such an induction.  The stated uniqueness hypothesis therefore
implies $W^\gamma=gW$ for some $g\in G$.
But
\[
 \operatorname{Stab}_G(W^\gamma)=H,
 \qquad\operatorname{Stab}_G(gW)=gHg^{-1},
\]
so $g\in N_G(H)$.  The two representations of $H$ on this
common block are respectively $\tau^\gamma$ and a conjugate of
$\tau$ by $g$ (or by its inverse, according to the induction
convention).  Replacing $g$ by its inverse if needed establishes
(BI.51).  This completes both sufficient-condition arguments.
\end{proof}

\begin{corollary}[Necessary asymmetry of the remaining twelve-branch cases]
For a rank-nine source with a strictly larger inducing field and
a twelve-branch rational intermediate base, the inducing subgroup
is nonnormal and its specified system of imprimitivity cannot be
unique among systems of that block dimension.  For a primitive
inducing block it cannot be the unique primitive system either.
This leaves the nonnormal, nonunique cases open; it does not
exclude them by a field-descent assumption.
\end{corollary}
